% ---
% titre: Problème de nombres relatifs - les passagers du bus
% theme: NO
% chapitre: Nombres relatifs
% sous_chapitre: Addition et soustraction
% niveau: [10e]
% type: EVAL
% tags: [c-addition,c-soustraction,c-nombres-relatifs,p-question-TS]
% difficulte: 1
% corrige: true
% ---

\question[1\half]
Sept personnes se sont installées dans le bus au moment où il démarre.

\begin{itemize}
\item Au premier arrêt, 3 personnes descendent du bus et 6 personnes montent.
\item Au deuxième arrêt, personne ne monte mais 4 personnes descendent.
\item Au troisième arrêt, 9 personnes montent et 6 descendent.
\item Au terminus, tout le monde descend !
\end{itemize}

\vspace{10pt}
 Combien de personnes descendent au terminus ?

\vspace{10pt}{\footnotesize\raggedleft Espace pour ta démarche et tes calculs\par}
\begin{solutionorgrid}[80pt]
7 - 3 + 6 - 4 + 9 - 6   = 9 \textbf{(0.5pt pour le calcul, 0.5pt pour la réponse)}
\end{solutionorgrid}

\begin{tcolorbox}[enhanced jigsaw, interior hidden, colframe=box, parbox=false]%
Réponse: 
\sol{\vspace{20pt}}{9 personnes descendent au terminus. \textbf{(0.5pt)}}
\end{tcolorbox}
